
\section{Related Work}
\label{sec:related}
\noindent 
Related work comprises the fields of 2.5D software maps and metaphor-based visualization of software system evolution.

\paragraph*{Interactive Software Maps}

Treemaps are commonly used to visualize hierarchical datasets. 
The original space-filling algorithm introduced by \cite{Shneiderman1992} slices parent nodes recursively by their child nodes' properties, alternating in horizontal and vertical direction. 
Due to undesirable aspect ratios produced by this slice and dice algorithm, Bruls \etal as well as Bederson and Shneiderman proposed squarified treemaps \cite{Bruls1999} and strip treemaps \cite{Bederson2002}. 
Furthermore, while preserving the visualization of the hierarchy, Fekete and Plaisant added visual attributes, such as color and size, to treemap items \cite{Fekete2002}. 
The first 2.5D treemap was described in the \emph{Step Tree} approach presented by Bladh \etal \cite{Bladh2004}. 
The geometric representation of sub-directories is stacked on top of their parent directories. 
This extension to the third dimension enables metaphor-based visualizations as shown in \emph{Information Pyramids} \cite{Andrews1997} and \emph{Code City} \cite{Wettel2008}, a city metaphor for 3D treemaps in which the height of the treemap's cuboids is used to map an additional software system metric. 
Similar to actual cities, these code cities are partitioned into several districts, including downtown areas and suburbs.
Referenced as \emph{software maps}, the city metaphor was adopted and extended by several other publications \cite{Liggesmeyer2009,Bohnet2011}.
Bohnet and D{\"o}llner describe the usage of software maps as a decision-making tool.
Recently, a comparative performance evaluation of rendering techniques for 2.5D treemaps was presented \cite{Trapp2013}.
With it, a deferred shape-generation algorithm outperforming existing approaches is presented.
The basic idea is to generate and render an attributed point-cloud and to expand the points to cuboids on the GPU.

\paragraph*{Visualization of Software Evolution}

The visualization of software evolution implies an additional time dimension.
Most techniques visualize the change process by animating the transition phase between two software revisions.
The \emph{VERSO visualization tool} \cite{Langelier2008} uses linear interpolation of color, height, and twist between two revisions. 
Another approach uses an Evolution Storyboard consisting of an ordered sequence of animated panels showing the structural changes during a time period \cite{Beyer2006}.
In contrast to animation, the Evolution Matrix approach displays the entire evolution of the software system in a single image \cite{Lanza2001}. 
Rectangles are used to map the number of methods to width and the number of attributes to height. 
Each revision is visualized next to its consecutive revision in a separate column.

\paragraph*{Metaphors in Software Visualization}

Besides the city metaphor, other metaphors for software system visualization were proposed. 
Balzer \etal presented a software landscape metaphor \cite{Balzer2004} using landscape-like distributions of three-dimensional objects on a 2D reference plane. 
Depending relationships between subsystems are visualized using a hierarchical network. 
Further, Graham \etal proposed a solar-system metaphor \cite{Graham2004} for object-oriented software metrics.
Planets represent classes, orbits depict the inheritance level within a package: planets in each orbit represent an inheritance hierarchy, while the lines-of-code of a class are mapped to the planet size. 
A thoroughly written survey of visualization techniques for static and evolutionary information of software systems is presented by Caserta and Zendra \cite{Caserta2011}.
Holten \etal presented a 2D treemap visualization of software metrics using the surface structure of treemap elements by applying texture- and bump mapping to them \cite{Holten2005}. 
Their approach is based on the fact that the human visual system can rapidly process visual cues such as shading and textures \cite{Enns1990,Healey1998,Healey1999}.
Panas \etal use a realistic city metaphor to map code quality metrics to building textures \cite{Caserta2011}. 
Furthermore, they proposed visual effects such as fire and bolts to communicate hot spot information in code execution and frequent component modifications \cite{Panas2003}.
However, neither an implementation nor an evaluation of such visual mappings are provided.
An empirical study by Borgo \etal suggests that visual embellishments aid memorization performance in both accuracy and response time, though it might have a negative impact on speed of visual search \cite{Borgo2012}.
It further proves that visual embellishments aid concept grasping.

